% GENERATED FILE. Edit canonical lessons, then run npm run generate:books.
% canonical-source-hash: fnv1a64:7a94551cba551501
% canonical-lessons: JA-C305-yuubin, JA-C305-atesaki, JA-C305-sokutatsu, JA-C305-haitatsu, JA-C305-hisashiburi

\chapter{A Letter from Home}
\label{ch:ja-a-letter-from-home}
\hlchaptermodality{305}

\begin{chapteropening}
\textbf{By the end of this chapter:} \emph{I can say what came in the post and when.}

Complete the last lesson of chapter 305: I can say what came in the post and when.
\end{chapteropening}

\section[\texorpdfstring{\ja{ゆうびん} (yūbin)}{ゆうびん (yūbin)}]{\ja{ゆうびん} (yūbin) — the post; mail}

\label{lesson:JA-C305-yuubin}



\begin{quote}
Before the new one: say the Japanese for to peep, then the Japanese for needy.
\end{quote}

\subsection*{You'll want to know: \ja{ゆうびん}}
\textbf{\ja{ゆうびん}} — \emph{yūbin} — \textquotedblleft{}the post; mail\textquotedblright{}.

\begin{quote}
\textbf{\ja{けさ} \ja{ゆうびんが} \ja{きました。}} — \emph{kesa yūbin ga kimashita} — \textquotedblleft{}The post came this morning.\textquotedblright{}
\end{quote}

\begin{cousinweb}
\textbf{\ja{ゆうびん}} (\emph{yūbin}) is Sino-Japanese (kango): two Chinese characters read with their on'yomi, readings borrowed from Chinese. They are \emph{yū}, \textquotedblleft{}mail\textquotedblright{} (Middle Chinese \emph{hjuw}) and \emph{bin}, \textquotedblleft{}letter\textquotedblright{} (\emph{bjienH}) (KANJIDIC2; Middle Chinese in Baxter's transcription).
\end{cousinweb}

\begin{grammarlens}[title={What you've built}]
One more word for telling what happened.
\end{grammarlens}

\subsection*{Your turn}
\begin{itemize}
\raggedright
  \item \emph{Say it:} \emph{yūbin}
  \item \emph{Say it:} \emph{yūbin}, once more
  \item \emph{Say it:} \emph{kesa yūbin ga kimashita} — The post came this morning.
  \item \emph{Recall:} say the Japanese for to peep, then the Japanese for needy, then say \emph{yūbin} again
\end{itemize}

\begin{tcolorbox}[breakable,colback=teal!4,colframe=teal!35!black,title={Before you move on}]
What does \ja{ゆうびん} mean? (\textquotedblleft{}The post; mail\textquotedblright{}.) Say it once more.
\end{tcolorbox}
\section[\texorpdfstring{\ja{あてさき} (atesaki)}{あてさき (atesaki)}]{\ja{あてさき} (atesaki) — an address (on a letter)}

\label{lesson:JA-C305-atesaki}



\begin{quote}
Before the new one: say the Japanese for needy, then the Japanese for the post.
\end{quote}

\subsection*{You'll want to know: \ja{あてさき}}
\textbf{\ja{あてさき}} — \emph{atesaki} — \textquotedblleft{}an address (on a letter)\textquotedblright{}.

\begin{quote}
\textbf{\ja{あてさきを} \ja{まちがえました。}} — \emph{atesaki o machigaemashita} — \textquotedblleft{}I got the address wrong.\textquotedblright{}
\end{quote}

\begin{cousinweb}
\textbf{\ja{あてさき}} (\emph{atesaki}) is a native Japanese word (wago), inherited rather than borrowed from Chinese, built from \emph{ate}, \textquotedblleft{}aiming\textquotedblright{} (from \emph{ateru}), and \emph{saki}, \textquotedblleft{}the far end\textquotedblright{} (JMdict, KANJIDIC2).
\end{cousinweb}

\begin{grammarlens}[title={What you've built}]
One more word for telling what happened.
\end{grammarlens}

\subsection*{Your turn}
\begin{itemize}
\raggedright
  \item \emph{Say it:} \emph{atesaki}
  \item \emph{Say it:} \emph{atesaki}, once more
  \item \emph{Say it:} \emph{atesaki o machigaemashita} — I got the address wrong.
  \item \emph{Recall:} say the Japanese for needy, then the Japanese for the post, then say \emph{atesaki} again
\end{itemize}

\begin{tcolorbox}[breakable,colback=teal!4,colframe=teal!35!black,title={Before you move on}]
What does \ja{あてさき} mean? (\textquotedblleft{}An address (on a letter)\textquotedblright{}.) Say it once more.
\end{tcolorbox}
\section[\texorpdfstring{\ja{そくたつ} (sokutatsu)}{そくたつ (sokutatsu)}]{\ja{そくたつ} (sokutatsu) — express delivery}

\label{lesson:JA-C305-sokutatsu}



\begin{quote}
Before the new one: say the Japanese for the post, then the Japanese for an address (on a letter).
\end{quote}

\subsection*{You'll want to know: \ja{そくたつ}}
\textbf{\ja{そくたつ}} — \emph{sokutatsu} — \textquotedblleft{}express delivery\textquotedblright{}.

\begin{quote}
\textbf{\ja{この} \ja{てがみを} \ja{そくたつで} \ja{おくって} \ja{ください。}} — \emph{kono tegami o sokutatsu de okutte kudasai} — \textquotedblleft{}Please send this letter express.\textquotedblright{}
\end{quote}

\begin{cousinweb}
\textbf{\ja{そくたつ}} (\emph{sokutatsu}) is Sino-Japanese (kango): two Chinese characters read with their on'yomi, readings borrowed from Chinese. They are \emph{soku}, \textquotedblleft{}fast\textquotedblright{} (Middle Chinese \emph{suwk}) and \emph{tatsu}, \textquotedblleft{}arrive\textquotedblright{} (\emph{dat}) (KANJIDIC2; Middle Chinese in Baxter's transcription).
\end{cousinweb}

\begin{grammarlens}[title={What you've built}]
One more word for telling what happened.
\end{grammarlens}

\subsection*{Your turn}
\begin{itemize}
\raggedright
  \item \emph{Say it:} \emph{sokutatsu}
  \item \emph{Say it:} \emph{sokutatsu}, once more
  \item \emph{Say it:} \emph{kono tegami o sokutatsu de okutte kudasai} — Please send this letter express.
  \item \emph{Recall:} say the Japanese for the post, then the Japanese for an address (on a letter), then say \emph{sokutatsu} again
\end{itemize}

\begin{tcolorbox}[breakable,colback=teal!4,colframe=teal!35!black,title={Before you move on}]
What does \ja{そくたつ} mean? (\textquotedblleft{}Express delivery\textquotedblright{}.) Say it once more.
\end{tcolorbox}
\section[\texorpdfstring{\ja{はいたつ} (haitatsu)}{はいたつ (haitatsu)}]{\ja{はいたつ} (haitatsu) — a delivery}

\label{lesson:JA-C305-haitatsu}



\begin{quote}
Before the new one: say the Japanese for an address (on a letter), then the Japanese for express delivery.
\end{quote}

\subsection*{You'll want to know: \ja{はいたつ}}
\textbf{\ja{はいたつ}} — \emph{haitatsu} — \textquotedblleft{}a delivery\textquotedblright{}.

\begin{quote}
\textbf{\ja{ゆうびんの} \ja{はいたつは} \ja{ごぜんです。}} — \emph{yūbin no haitatsu wa gozen desu} — \textquotedblleft{}The post is delivered in the morning.\textquotedblright{}
\end{quote}

\begin{cousinweb}
\textbf{\ja{はいたつ}} (\emph{haitatsu}) is Sino-Japanese (kango): two Chinese characters read with their on'yomi, readings borrowed from Chinese. They are \emph{hai}, \textquotedblleft{}distribute\textquotedblright{} (Middle Chinese \emph{phwojH}) and \emph{tatsu}, \textquotedblleft{}arrive\textquotedblright{} (\emph{dat}) (KANJIDIC2; Middle Chinese in Baxter's transcription).
\end{cousinweb}

\begin{grammarlens}[title={What you've built}]
One more word for telling what happened.
\end{grammarlens}

\subsection*{Your turn}
\begin{itemize}
\raggedright
  \item \emph{Say it:} \emph{haitatsu}
  \item \emph{Say it:} \emph{haitatsu}, once more
  \item \emph{Say it:} \emph{yūbin no haitatsu wa gozen desu} — The post is delivered in the morning.
  \item \emph{Recall:} say the Japanese for an address (on a letter), then the Japanese for express delivery, then say \emph{haitatsu} again
\end{itemize}

\begin{tcolorbox}[breakable,colback=teal!4,colframe=teal!35!black,title={Before you move on}]
What does \ja{はいたつ} mean? (\textquotedblleft{}A delivery\textquotedblright{}.) Say it once more.
\end{tcolorbox}
\section[\texorpdfstring{\ja{ひさしぶり} (hisashiburi)}{ひさしぶり (hisashiburi)}]{\ja{ひさしぶり} (hisashiburi) — after a long time; for the first time in ages}

\label{lesson:JA-C305-hisashiburi}



\begin{quote}
Before the new one: say the Japanese for express delivery, then the Japanese for a delivery.
\end{quote}

\subsection*{You'll want to know: \ja{ひさしぶり}}
\textbf{\ja{ひさしぶり}} — \emph{hisashiburi} — \textquotedblleft{}after a long time; for the first time in ages\textquotedblright{}.

\begin{quote}
\textbf{\ja{ともだちから} \ja{ひさしぶりの} \ja{てがみが} \ja{きました。}} — \emph{tomodachi kara hisashiburi no tegami ga kimashita} — \textquotedblleft{}A letter came from my friend, the first in ages.\textquotedblright{}
\end{quote}

A na-adjective: \ja{な} before a noun.

\begin{cousinweb}
\textbf{\ja{ひさしぶり}} (\emph{hisashiburi}) is a native Japanese word (wago), inherited rather than borrowed from Chinese, built from \emph{hisashi}, \textquotedblleft{}long-lasting\textquotedblright{} (from \emph{hisashii}), and \emph{buri}, \textquotedblleft{}after an interval\textquotedblright{} (JMdict, KANJIDIC2).
\end{cousinweb}

\begin{grammarlens}[title={What you've built}]
One more word for telling what happened.
\end{grammarlens}

\subsection*{Your turn}
\begin{itemize}
\raggedright
  \item \emph{Say it:} \emph{hisashiburi}
  \item \emph{Say it:} \emph{hisashiburi}, once more
  \item \emph{Say it:} \emph{tomodachi kara hisashiburi no tegami ga kimashita} — A letter came from my friend, the first in ages.
  \item \emph{Recall:} say the Japanese for express delivery, then the Japanese for a delivery, then say \emph{hisashiburi} again
\end{itemize}

\begin{tcolorbox}[breakable,colback=teal!4,colframe=teal!35!black,title={Before you move on}]
What does \ja{ひさしぶり} mean? (\textquotedblleft{}After a long time; for the first time in ages\textquotedblright{}.) Say it once more.
\end{tcolorbox}
